\documentclass[letterpaper,10pt,notitlepage,fleqn]{article}

%\usepackage{nopageno} %gets rid of page numbers
\usepackage{alltt}                                           
\usepackage{float}
\usepackage{color}
\usepackage{indentfirst}
\usepackage{url}
\usepackage{balance}
\usepackage[TABBOTCAP, tight]{subfigure}
\usepackage{enumitem}
\usepackage{pstricks, pst-node}
\usepackage{geometry}
\geometry{textheight=9in, textwidth=6.5in} %sets 1" margins 
\newcommand{\cred}[1]{{\color{red}#1}} %command to change font to red
\newcommand{\cblue}[1]{{\color{blue}#1}} % ...blue
\usepackage{hyperref}
\usepackage{textcomp}
\usepackage{listings}
\usepackage{graphicx}

\def\name{Sam Quinn}

\parindent = 0.4444 in
\parskip = 0.2 in

\begin{document}
\begin{titlepage}
\vspace*{\fill}

\newcommand{\HRule}{\rule{\linewidth}{0.5mm}} % Defines a new command for the horizontal lines, change thickness here

\center % Center everything on the page

%----------------------------------------------------------------------------------------
%TITLE SECTI   ON
%----------------------------------------------------------------------------------------

\includegraphics[scale=.5]{pcc.eps}
\HRule \\[0.4cm]
{ \huge \bfseries Precision Castparts Corp. Report}\\[0.4cm] % Title of your document

%----------------------------------------------------------------------------------------
%HEADING SECTIONS
%----------------------------------------------------------------------------------------

\textsc{\LARGE Sam Quinn}\\[0.5cm] % Name of your university/college
\textsc{\Large MECOP Internship \#2}\\[0.5cm] % Major heading such as course name
\textsc{\large Summer/Fall 2015}\\[0.5cm] % Minor heading such as course title


\HRule \\[1.5cm]
%----------------------------------------------------------------------------------------
%AUTHOR SECTION
%------------------------------------ ----------------------------------------------------

\begin{minipage}{0.4\textwidth}
\begin{flushleft} \large
\emph{Intern:}\\
        \textbf{Sam \textsc{Quinn}} \\ % Your name
        {\small Computer Science}
        \end{flushleft}
        \end{minipage}
        ~
        \begin{minipage}{0.4\textwidth}
        \begin{flushright} \large
        \emph{Mentors:} \\
            \textbf{Steve \textsc{Jessie}} \\ % Supervisor's Name
            {\small sjessie@pccstructurals.com \\
            (503)777-8811} \\
            
            \textbf{Dave \textsc{Rajkumar}}
            {\small drajkumar@pccstructurals.com \\
            (503)777-8816}
            \end{flushright}
            \end{minipage}\\[3cm]

            % If you don't want a supervisor, uncomment the two lines below and remove the section above
            %\Large \emph{Author:}\\
                %John \textsc{Smith}\\[3cm] % Your name

                %----------------------------------------------------------------------------------------
                %DATE SECTION
                %-----------------    -----------------------------------------------------------------------

{\large \today}\\[3cm] % Date, change the \today to a set date if you want to be precise

%----------------------------------------------------------------------------------------
%LOGO SECTION
%------   ----------------------------------------------------------------------------------

%\includegraphics{Logo}\\[1cm] % Include a department/university logo - this will require the graphicx package

%----------------------------------------------------------------------------------------

\vfill % Fill the rest of the page with whitespace



\end{titlepage}

\tableofcontents
\newpage

\section{Introduction}
\indent For my second MECOP internship I was fortunate to be placed at Precision Castpart Corp. This 
was my second time being immersed into a professional atmosphere in the field that 
I am studying and I enjoyed it equally as the first. I am computer science student with a focus 
in security and even thought that Precision Castparts Corp. was more of an old school 
non-tech company, I was greatly mistaken. Precision Castparts Corp. had a lot of advanced 
technology and very intricate security, which was great for me to get hands on with. It was 
good for me to learn the basics of how a large corporation like Precision Castpart Corp. 
handles security. I now feel more well rounded and have at least some experience relating to
my security focus.

\section{Precision Castparts Corp.}

\subsection{What does Precision Castpart Corp. do?}
\indent When I was first placed at Precision Castpart Corp. I had not heard much about them. I 
definitely was not a ware that they were one of Oregon's two fortune 500 companies.
Precision Castparts Corporation specializes in the casting of metal aerospace parts. 
They are the world leader in investment casting, forged components, and airfoil castings. 
The company was started over 60 years ago in Portland, OR by Joseph Cox, the owner of 
Oregon Saw Chain. The initial intention was to cast a special cutter for his chain saw 
products. The company began focusing their time on developing the casting operations 
and looking for more outside work. This portion of the company became so successful that 
it branched out and became a separate business called Precision Castparts Co. The company 
began pushing to create larger castings and, by 1962, was able to pour a casting weighing 
1,000 pounds. PCC soared above all of the competition, and later signed contracts with 
General Electric, Boeing, Pratt \& Whitney, and several other aerospace companies. The company 
began developing their processes with titanium, nickel, and stainless steel alloys. PCC soon
became the world leader of manufacturing large titanium castings as well. The company currently 
has several campuses throughout Portland, and has recently been acquired by Warren 
Buffett's Berkshire Hathaway.

\subsection{MIS}
\indent During my internship at Precision Castparts Corp I was exposed to many 
different departments. The casting process has not changed much over the 60 years 
Precision Castparts Corp has been around, but what has changed a lot has been the 
technology that they have been using to refine the casting process and this is where 
the MIS department comes in. All of the technology is managed by the MIS 
(Management Information System) department. For a large 
portion of my internship I was stationed at the helpdesk. I liked the atmosphere 
of the helpdesk and the opportunity to get hands on experience with technology all
around the plant. I would help individual users with minor problems like email 
to large issues that were halting production like robotic arms. 

\subsection{Mentor/Supervisors}
\indent My mentors were the managers of the MIS department Steve Jessie and Dave 
Rajkumar. Both Steve and 
Dave were responsible for keeping all of Precision Castparts Corp technology up 
to date and secure. There management encompassed the LAN Admins, Programmers, 
Helpdesk, and the Telecom employees. Even though I was often at my desk in the 
helpdesk which was separate from the MIS department we kept a close relationship 
with the other programmers, LAN admins, and techs. I worked with a lot of individuals 
during my time at Precision Castpart Corp. Often times when a problem would arise 
we would all have a role in getting the issue resolved. 

\section{Projects}

\subsection{Aras Innovator Upgrade}
\indent The first large project that I had the opportunity to work on was upgrading 
Aras Innovator. The drive for the MIS department to upgrade Innovator was that the 
current version when I began at PCC only worked on Internet Explorer 9. We wanted 
to move all of our computers to Internet Explorer 11 and to do this we had to make 
sure that Innovator would continue to work. Innovator is a unique process control 
visualization tool that is used company wide and a seamless upgrade with no errors
was required. I worked very closely with Kyle Grotjohn who was the lead in the 
upgrade. We both had to create a comprehensive test plan that would get as close 
to 100\% coverage as possible. This plan was huge we if there was a check box that 
could be clicked there was a whole test just to test that one check box. I spent 
about 4 weeks going through the test plan and creating bug reports for any errors 
I came across. The bug report contained screenshots, the steps taken to reproduce 
the error, and my description of the what was supposed to happen. Aras corp were 
actually the ones in charge of fixing the bugs, so I would package all of my 
reports up and send it to them. 
\\
\indent This was the first time I have been introduced to the Aras Innovator Application. 
Since Kyle had begun to test this tool before me, he was a great mentor on the application. 
It was really nice to have some tests in place before I started, and to view how a professional
laid out the base of the test plan. Mimicking the style helped me create and execute 
a well designed test plan.
\\
\indent One of the biggest takeaways from this project was the importance of clear 
and easy to read documentation. In my past projects that I developed tests for I 
never created a test plan which probably resulted in many functions not getting coverage.
Studying the test plan before beginning my work on this project, I noticed that 
there were clear comments everywhere. Since I mimicked Kyle's testing style, I too 
used an abundance of comments. I soon came to the realization of how much easier 
it is to understand and take care of bugs when you know where they come from. \\


\textbf{Things Learned:}
\begin{itemize} 
        \item .Net framework
        \item Internet Explorer compatibility
        \item Creating Bug reports
        \item Work with others following a plan
    \end{itemize}

\subsection{Security Audit}
\indent With many of the contracts that PCC has the cyber security has to be top 
notch.  PCC would conduct bimonthly security audits that would check every computer 
that is in our domain. We would get separate lists of computers that have reported 
computers that are missing our mandatory security software. Each list would have 
a lot of computers that would require us to manually connect to each machine and 
verify that the computer either reported a false-positive or was indeed missing 
the software. I saw this as a great opportunity to implement automation. I began 
creating scripts that would make the audit process quicker and more accurate. The 
first Script that I wrote would take all the computer names and try to gather more 
info on them. The script would ping each computer to see if they are currently 
online, get the last user who logged into the workstation, and the last time time 
the computer was logged into. 
\\
\indent
At first the other PC techs were skeptical to use my scripts since they have been 
doing the security audits many times in the past and did not want to cause more 
problems, but after a while they began to see the usefulness in what they did. 
Each script needed to increase the productivity of the audit process. I spent time 
developing the scripts which put me behind schedule on the first audit but the 
audits that came prior both the scripts increase my productivity tremendously.
\\
\indent Before the first security audit I had not used Powershell to create scripts.
It was a slight learning curve that I had to over come. I am very familiar with the 
Unix shell environment and Powershell shared many of the same characteristics. One 
challenge that I had to overcome was that Windows would not actually run any Powershell 
script that had not been signed. This was a problem since I did not have access to 
a Windows development key for signing, I ended up creating a wrapper batch script 
that would enable unsigned execution and then after executing set the execution 
policy back to default. Needless to say writing the batch script wrapper was more 
difficult than the Powershell script. The Batch script syntax was very different 
from the familiar Unix shell environment. 

\textbf{Things Learned:}
\begin{itemize} 
        \item Powershell scripts
        \item SCCM reporting
        \item Active Directory
        \item Windows CLI
    \end{itemize}

\section{Minor Projects}

\subsection{TWiki Upgrade}
\indent During my down time while working with Sam debugging the Fill Wrap feature I had 
the opportunity to gain knowledge in website hosting. I was asked to investigate 
what needed to be done to upgrade DFM QA's internal Wiki site. The DFM QA Wiki 
site was created with TWiki and open source wiki interface and hosted locally at Mentor 
Graphics. Before diving head first into this project, I had no prior experience working 
with Apache, CGI scripts, or really any web tools for that matter, but I was very 
excited to learn. \\
\indent I spent about a week creating an exact replica of the DFM QA TWiki site on my local 
machine with the updated version. I ran into a lot of problems with permissions 
and how Apache can own files that even root can't access. With help from Sam and 
ICV BIG, I eventually got a working copy running on my computer. Since the TWiki 
was accessed by all of the DFM QA team, there was too much information for me to 
verify everything copied correctly during the upgrade. The way Sam and I tested 
everything was correct, was we did a live switch and allowed all of the team to 
poke around the new version to check that all of their paths and normal day links 
continued to work. After we got our good results back from the team, we copied my 
local copy to the permanent server location. \\

\textbf{Things Learned:}
    \begin{itemize} 
        \item Apache configuration
        \item Unix Permissions
        \item CGI scripts
        \item General web hosting 
    \end{itemize}

\subsection{Test Suite Search Tool}
\indent After working with Dave in my Tinstall project, Dave approached me with a problem 
that he thought I could help with. Dave's problem was that when he finds a bug within 
the test suite he often times did a series of greps and pipes to narrow down his 
search. Since there wasn't really any built-in UNIX or inside tool that could search 
for multiple terms, I decided to help him with this. \\ 
\indent Dave gave me the freedom to use whatever language I thought best, since I was writing 
this program from scratch. After enjoying learning TCL so much in the Tinstall 
program, I decided to learn a new language and went with Perl for the code and Tk 
for the GUI. Dave thought Perl and Tk were a good choice and after explaining to 
him that I have never programmed in Perl before he showed me our company library, 
which was a great resource in later projects.

\textbf{Things Learned:}
\begin{itemize} 
        \item Tk
        \item Perl
        \item Making a GUI
        \item Grep to a higher degree
    \end{itemize}

\section{Conclusion}
\subsection{What I learned}
\indent During my 6 month internship I learned a comparable amount of degree related knowledge
than I would have at school. The main difference was that the things I learned while 
Working at Mentor Graphics no class could teach. One of my most exciting things 
I learned during my internship is how to be efficient. I consider myself 100\% more 
comfortable in a UNIX environment. I believe that I will be a much better programmer 
and will be able to produce a higher quality of work from the small tips and tricks 
that were taught to me during my time at Mentor Graphics. The real world work environment seems 
far less scary to me now that I have had a slight taste of it. It was very neat 
to see how my background knowledge from school got integrated into the work place.
I would often times recall information from previous classes that would help me 
solve problems I faced during my internship. 
\\
\subsection{How the internship benefited me}
\indent This internship has greatly benefited me in many ways. One of the greatest 
benefits I will takeaway from this internship is knowing that I can work in the 
industry. It was great to see that I was working on real projects that full time 
employees were working on and helped solve unsolved problems. I liked working on 
nonstandard problems in contrast to how at school I felt like I was finishing 
a homework assignment that I know every other CS student has solved before me. 
\\
\subsection{How I benefited Mentor Graphics}
\indent I really didn't know how important my contributions were to Mentor Graphics until 
I was talking to  one of the employees about what I was working on with the new Wrap function and he said "You 
are working on probably the most important update in the upcoming release". After 
I heard this I truly felt like Mentor Graphics valued my work and that I was actually
making a difference. I created more than 30 test cases, upgraded a site most people 
on the QA team access daily, and helped find which functions that need to be tested.\\ 
\subsection{Presentations}
\indent After finishing Mentor Graphic's internal tool, Tinstall, I was asked to present to 
a majority of the QA employees explaining my tool. This presentation took place via 
telecommunication with people viewing from Mentor Graphics locations all over the 
world.
\end{document}

